\documentclass[12pt,a4paper]{article}
\usepackage[utf8]{inputenc}
\usepackage[T1]{fontenc}
\usepackage{geometry}
\usepackage{graphicx}
\usepackage{hyperref}
\usepackage{enumitem}
\usepackage{xcolor}
\usepackage{tcolorbox}
\usepackage{booktabs}
\usepackage{longtable}
\usepackage{fancyhdr}
\usepackage{titlesec}
\usepackage{amsmath}
\usepackage{listings}

% Page geometry
\geometry{margin=1in}

% Colors
\definecolor{algovizpurple}{RGB}{139, 92, 246}
\definecolor{algovizorange}{RGB}{245, 158, 11}
\definecolor{algovizgreen}{RGB}{16, 185, 129}
\definecolor{algovizblue}{RGB}{59, 130, 246}
\definecolor{algovizred}{RGB}{239, 68, 68}
\definecolor{darkbg}{RGB}{14, 17, 23}
\definecolor{lighttext}{RGB}{203, 213, 225}

% Hyperlink setup
\hypersetup{
    colorlinks=true,
    linkcolor=algovizpurple,
    urlcolor=algovizblue
}

% Header/Footer
\pagestyle{fancy}
\fancyhf{}
\fancyhead[L]{\textcolor{algovizpurple}{\textbf{AlgoViz}} - Presentation Strategy}
\fancyhead[R]{\thepage}
\fancyfoot[C]{\small See the Signal in the Noise}

% Title format
\titleformat{\section}
{\color{algovizpurple}\Large\bfseries}
{\thesection}{1em}{}

\titleformat{\subsection}
{\color{algovizblue}\large\bfseries}
{\thesubsection}{1em}{}

\titleformat{\subsubsection}
{\color{algovizorange}\normalsize\bfseries}
{\thesubsubsection}{1em}{}

% Tip box
\newtcolorbox{tipbox}[1][]{
    colback=algovizgreen!10,
    colframe=algovizgreen,
    fonttitle=\bfseries,
    title=#1,
    arc=3mm,
}

% Warning box
\newtcolorbox{warningbox}[1][]{
    colback=algovizorange!10,
    colframe=algovizorange,
    fonttitle=\bfseries,
    title=#1,
    arc=3mm,
}

% Important box
\newtcolorbox{importantbox}[1][]{
    colback=algovizpurple!10,
    colframe=algovizpurple,
    fonttitle=\bfseries,
    title=#1,
    arc=3mm,
}

\begin{document}

%===============================================================================
% TITLE PAGE
%===============================================================================
\begin{titlepage}
    \centering
    \vspace*{2cm}
    
    {\Huge\bfseries\textcolor{algovizpurple}{AlgoViz}\par}
    \vspace{0.5cm}
    {\Large\textcolor{lighttext}{See the Signal in the Noise}\par}
    
    \vspace{2cm}
    
    {\LARGE\bfseries Presentation Strategy Guide\par}
    \vspace{0.5cm}
    {\large Complete Reference for Dashboard Demonstration\par}
    
    \vspace{3cm}
    
    {\large\textbf{Created for:} DVA Final Project Presentation\par}
    \vspace{0.3cm}
    {\large\textbf{Team:} Kartik Joshi, Aditya Chitale, Krishna Patel\par}
    
    \vspace{2cm}
    
    \begin{importantbox}[\textbf{Purpose of This Document}]
    This document provides a complete, step-by-step guide to understand and present the AlgoViz dashboard. It covers every single feature, filter, chart, KPI, and technical component in simple, easy-to-understand language. Use this as your presentation script and reference guide.
    \end{importantbox}
    
    \vfill
    {\small February 2026\par}
\end{titlepage}

%===============================================================================
% TABLE OF CONTENTS
%===============================================================================
\tableofcontents
\newpage

%===============================================================================
% SECTION 1: INTRODUCTION - WHAT IS ALGOVIZ?
%===============================================================================
\section{Introduction: What is AlgoViz?}

\subsection{The Big Picture}
AlgoViz is a \textbf{professional-grade algorithmic trading visualization dashboard} that helps traders and analysts understand what's happening in financial markets in real-time. 

\begin{tipbox}[\textbf{Simple Explanation}]
Think of AlgoViz like a car's dashboard, but for the stock/crypto market. Just like a car dashboard shows you speed, fuel level, and engine temperature, AlgoViz shows you price movements, trading activity, and market conditions -- all updating in real-time.
\end{tipbox}

\subsection{What Makes AlgoViz Unique?}

Unlike regular trading dashboards, AlgoViz has several \textbf{key differentiators}:

\begin{enumerate}[label=\arabic*.]
    \item \textbf{Deep Learning Integration:} Uses LSTM neural networks with attention mechanism to predict price direction
    \item \textbf{Real-Time WebSocket Connection:} Live data streaming directly from Binance exchange
    \item \textbf{Trading Intelligence Panel:} AI-generated actionable insights, not just raw data
    \item \textbf{Strategy Backtester:} Test trading strategies before risking real money
    \item \textbf{Two Modes:} Works with both demo data (for learning) and live market data
    \item \textbf{Professional Design:} Dark theme with glassmorphism effects, suitable for professional trading environments
\end{enumerate}

\subsection{Technology Stack}

\begin{table}[h]
\centering
\begin{tabular}{ll}
\toprule
\textbf{Component} & \textbf{Technology Used} \\
\midrule
Frontend Framework & Streamlit (Python-based web framework) \\
Charts & Plotly (Interactive JavaScript charts) \\
Backend & Python 3.11 \\
Data Source & Binance WebSocket API \\
Deep Learning & Custom LSTM implementation in NumPy \\
Containerization & Docker \& Docker Compose \\
\bottomrule
\end{tabular}
\caption{Technology Stack Overview}
\end{table}

%===============================================================================
% SECTION 2: HOW TO ACCESS THE DASHBOARD
%===============================================================================
\section{How to Access the Dashboard}

\subsection{Using Docker (Recommended)}
\begin{enumerate}
    \item Open terminal in the \texttt{hft-dashboard} folder
    \item Run: \texttt{docker-compose up -d}
    \item Open browser and go to: \texttt{http://localhost:8501}
\end{enumerate}

\subsection{Running Locally}
\begin{enumerate}
    \item Activate virtual environment: \texttt{.venv\textbackslash Scripts\textbackslash activate}
    \item Navigate to src folder: \texttt{cd hft-dashboard/src}
    \item Run: \texttt{streamlit run app.py}
    \item Dashboard opens automatically in browser
\end{enumerate}

\begin{warningbox}[\textbf{For Live Data Mode}]
You need an active internet connection. The dashboard connects to Binance servers to get real-time Bitcoin price data.
\end{warningbox}

%===============================================================================
% SECTION 3: SIDEBAR - NAVIGATION AND FILTERS
%===============================================================================
\section{Sidebar: Navigation and Filters}

The sidebar on the left side is your \textbf{control center}. Let's understand each part:

\subsection{Data Mode Selection}

This is the most important choice you make. There are \textbf{two modes}:

\subsubsection{Static Demo Mode}
\begin{itemize}
    \item Uses \textbf{synthetic (fake) data} generated by the dashboard
    \item Good for \textbf{learning and demonstration}
    \item Does \textbf{NOT require internet}
    \item Shows how the dashboard works without needing real market data
    \item Has a ``Data Quality'' tab in Analytics (not available in Live mode)
\end{itemize}

\subsubsection{Live Data Mode}
\begin{itemize}
    \item Connects to \textbf{Binance cryptocurrency exchange}
    \item Shows \textbf{actual real-time} Bitcoin (BTC/USDT) trading data
    \item Requires \textbf{active internet connection}
    \item Data updates every \textbf{250-500 milliseconds}
    \item Has a ``Live Feed'' page (not available in Static mode)
\end{itemize}

\begin{importantbox}[\textbf{How to Switch Modes}]
Look in the sidebar under ``Data Mode''. Select either:
\begin{itemize}
    \item ``\textbf{📊 Static Demo}'' - For safe demonstration
    \item ``\textbf{📡 Live Data}'' - For real market data
\end{itemize}
\end{importantbox}

\subsection{Market Scenario (Static Mode Only)}

When in Static Demo mode, you can choose different market scenarios:

\begin{table}[h]
\centering
\begin{tabular}{lp{10cm}}
\toprule
\textbf{Scenario} & \textbf{What It Simulates} \\
\midrule
Normal & Typical market conditions with moderate price movements \\
Volatile & High volatility - prices swing up and down rapidly \\
Trending Up & Bull market - prices generally moving upward \\
Trending Down & Bear market - prices generally moving downward \\
Low Liquidity & Thin market - wide spreads, fewer traders \\
\bottomrule
\end{tabular}
\caption{Available Market Scenarios}
\end{table}

\subsection{Page Navigation}

The sidebar shows different pages you can navigate to:

\begin{table}[h]
\centering
\begin{tabular}{lp{8cm}l}
\toprule
\textbf{Page} & \textbf{Purpose} & \textbf{Available In} \\
\midrule
🏠 Home & Landing page with features overview & Both modes \\
📊 Dashboard & Main trading view with KPIs and charts & Both modes \\
📡 Live Feed & Raw trade/order book stream & Live mode only \\
📈 Analytics & Advanced analysis and backtesting & Both modes \\
⚙️ Settings & Configuration options & Both modes \\
\bottomrule
\end{tabular}
\caption{Page Navigation Overview}
\end{table}

%===============================================================================
% SECTION 4: HOME PAGE
%===============================================================================
\section{Home Page}

The Home page is the \textbf{landing page} that introduces AlgoViz.

\subsection{What You'll See}

\begin{enumerate}
    \item \textbf{Hero Section:} Large animated logo with tagline ``See the Signal in the Noise''
    \item \textbf{Ticker Bar:} Scrolling text showing key features (BTC/USDT, LSTM, 250ms refresh)
    \item \textbf{Statistics Cards:} 
    \begin{itemize}
        \item \textbf{10+} Chart Types
        \item \textbf{2} Deep Learning Models
        \item \textbf{4} Trading Strategies
        \item \textbf{$\infty$} Insights (unlimited AI-generated insights)
    \end{itemize}
    \item \textbf{Feature Cards:} Six cards explaining main features
    \item \textbf{Call-to-Action:} Instructions to explore the dashboard
    \item \textbf{Team Credits:} Names of team members
\end{enumerate}

\begin{tipbox}[\textbf{Presentation Tip}]
Use the Home page as your introduction. Show the animated elements and feature cards to give an overview before diving into the actual dashboard.
\end{tipbox}

%===============================================================================
% SECTION 5: DASHBOARD PAGE - THE MAIN VIEW
%===============================================================================
\section{Dashboard Page: The Main Trading View}

This is the \textbf{heart of AlgoViz} where all the action happens.

\subsection{Header Section}

At the very top, you'll see:
\begin{itemize}
    \item \textbf{AlgoViz Logo:} With ``See the Signal in the Noise'' tagline
    \item \textbf{Data Source:} Shows ``Binance WebSocket''
    \item \textbf{Live Badge:} Green dot showing ``LIVE'' if connected, red if offline
\end{itemize}

\subsection{Connection Status Banner}
Shows:
\begin{itemize}
    \item Connection status (Connected/Reconnecting/Disconnected)
    \item Last update timestamp (shows exact time data was received)
    \item Date in YYYY-MM-DD format
\end{itemize}

\subsection{Chart Controls (Filters)}

Below the header, there are \textbf{four filter dropdowns}:

\subsubsection{Price Chart Window}
Controls how much history the Price \& VWAP chart shows:
\begin{itemize}
    \item \textbf{30 sec} - Last 30 seconds of data
    \item \textbf{1 min} - Last 1 minute
    \item \textbf{2 min} - Last 2 minutes (default)
    \item \textbf{5 min} - Last 5 minutes
\end{itemize}

\subsubsection{Volatility Window}
Controls the volatility chart time range:
\begin{itemize}
    \item \textbf{1 min} - 1 minute rolling window
    \item \textbf{2 min} - 2 minute rolling window (default)
    \item \textbf{5 min} - 5 minute rolling window
    \item \textbf{10 min} - 10 minute rolling window
\end{itemize}

\subsubsection{Line Style}
Controls how lines are drawn on charts:
\begin{itemize}
    \item \textbf{Smooth (Spline)} - Curved lines, looks prettier (default)
    \item \textbf{Sharp (Linear)} - Straight lines between points
    \item \textbf{Step} - Staircase pattern, shows exact changes
\end{itemize}

\subsubsection{Fullscreen-Safe Updates}
A status indicator showing that charts update without interrupting fullscreen view.

%===============================================================================
% SECTION 6: KPI CARDS (Key Performance Indicators)
%===============================================================================
\section{KPI Cards: Understanding the Numbers}

The dashboard shows \textbf{8 KPI cards} arranged in two rows.

\subsection{Row 1: Core Metrics}

\subsubsection{💰 PRICE (Gold accent)}
\begin{itemize}
    \item \textbf{What it shows:} Current Bitcoin price in USD
    \item \textbf{Format:} \$XX.XXK (e.g., \$87.35K means \$87,350)
    \item \textbf{Below it:} Shows ``BTC/USDT'' (the trading pair)
\end{itemize}

\begin{tipbox}[\textbf{How to Read}]
This is the current market price at which Bitcoin is trading. Updates every fraction of a second in live mode.
\end{tipbox}

\subsubsection{📊 CHANGE (Green/Red accent)}
\begin{itemize}
    \item \textbf{What it shows:} Price change percentage
    \item \textbf{Green with ▲:} Price went UP
    \item \textbf{Red with ▼:} Price went DOWN
    \item \textbf{Below it:} Actual dollar change (e.g., +\$15.50)
\end{itemize}

\subsubsection{📏 SPREAD (Purple accent)}
\begin{itemize}
    \item \textbf{What it shows:} Bid-Ask spread in basis points (bps)
    \item \textbf{1 bps = 0.01\%}
    \item \textbf{Status labels:}
    \begin{itemize}
        \item \textbf{Ultra-Tight} ($<$0.01 bps) - Exceptional liquidity
        \item \textbf{Tight} ($<$3 bps) - Good liquidity
        \item \textbf{Normal} (3-6 bps) - Average conditions
        \item \textbf{Wide} ($>$6 bps) - Poor liquidity, higher trading costs
    \end{itemize}
\end{itemize}

\begin{importantbox}[\textbf{What is Spread?}]
The spread is the difference between the highest price a buyer is willing to pay (bid) and the lowest price a seller is willing to accept (ask). Lower spread = better for traders because you lose less money when entering/exiting trades.
\end{importantbox}

\subsubsection{⚡ VELOCITY (Orange accent)}
\begin{itemize}
    \item \textbf{What it shows:} Trades per second
    \item \textbf{Format:} XX.X/s (e.g., 15.3/s means 15.3 trades per second)
    \item \textbf{Status labels:}
    \begin{itemize}
        \item \textbf{Normal} - Typical trading activity
        \item \textbf{High} - Above average activity (1.5x baseline)
        \item \textbf{Spike!} - Unusual surge (2x or more baseline)
    \end{itemize}
\end{itemize}

\subsection{Row 2: Advanced Metrics}

\subsubsection{📌 VWAP (Cyan accent)}
\begin{itemize}
    \item \textbf{What it shows:} Volume-Weighted Average Price
    \item \textbf{Status:} ``Above'', ``Below'', or ``Fair'' (relative to current price)
\end{itemize}

\begin{importantbox}[\textbf{What is VWAP?}]
VWAP is the average price weighted by trading volume. If you bought at VWAP, you paid the ``fair'' price that most volume traded at. 
\begin{itemize}
    \item \textbf{Price above VWAP:} You might be overpaying (overbought)
    \item \textbf{Price below VWAP:} You might be getting a deal (oversold)
\end{itemize}
\end{importantbox}

\subsubsection{⏱️ TWAP (Teal accent)}
\begin{itemize}
    \item \textbf{What it shows:} Time-Weighted Average Price
    \item \textbf{Status:} ``Above'', ``Below'', or ``Fair''
\end{itemize}

\begin{tipbox}[\textbf{VWAP vs TWAP}]
\textbf{VWAP} weights by volume (big trades matter more).\\
\textbf{TWAP} weights equally by time (every second matters the same).\\
Both help determine ``fair value'' of an asset.
\end{tipbox}

\subsubsection{⚖️ IMBALANCE (Pink accent)}
\begin{itemize}
    \item \textbf{What it shows:} Order book imbalance as percentage
    \item \textbf{Positive (+):} More buyers (bullish)
    \item \textbf{Negative (-):} More sellers (bearish)
    \item \textbf{Labels:} ``BUY'', ``SELL'', ``Buy'', ``Sell'' based on strength
\end{itemize}

\begin{importantbox}[\textbf{What is Order Book Imbalance?}]
Formula: $\text{Imbalance} = \frac{\text{Bid Volume} - \text{Ask Volume}}{\text{Bid Volume} + \text{Ask Volume}} \times 100\%$

\begin{itemize}
    \item \textbf{+50\% or more:} Strong buy pressure - price likely to go up
    \item \textbf{-50\% or less:} Strong sell pressure - price likely to go down
    \item \textbf{Near 0\%:} Balanced market
\end{itemize}
\end{importantbox}

\subsubsection{🔥 BUY PRESSURE (Indigo accent)}
\begin{itemize}
    \item \textbf{What it shows:} Percentage of volume that is buying
    \item \textbf{$>$55\%:} Buyers dominating (green, ``Buyers'')
    \item \textbf{$<$45\%:} Sellers dominating (red, ``Sellers'')
    \item \textbf{45-55\%:} Neutral market (purple, ``Neutral'')
\end{itemize}

%===============================================================================
% SECTION 7: DASHBOARD CHARTS
%===============================================================================
\section{Dashboard Charts: Visual Analytics}

The dashboard has \textbf{6 main charts} arranged in 3 rows of 2 columns.

\subsection{Chart 1: Price \& VWAP}

\subsubsection{What It Shows}
A line chart with two lines:
\begin{itemize}
    \item \textbf{White line:} Current price over time
    \item \textbf{Cyan line:} VWAP over time
    \item \textbf{Gradient fill:} Between the two lines, showing divergence
\end{itemize}

\subsubsection{Axes}
\begin{itemize}
    \item \textbf{X-axis:} Time (timestamps)
    \item \textbf{Y-axis:} Price in USD (automatically scaled)
\end{itemize}

\subsubsection{How to Interpret}
\begin{itemize}
    \item When \textbf{white line is above cyan}: Price is above VWAP = overbought
    \item When \textbf{white line is below cyan}: Price is below VWAP = potential value
    \item \textbf{Crossovers} (when lines cross): Potential momentum shift
\end{itemize}

\subsection{Chart 2: Bid-Ask Spread Timeline}

\subsubsection{What It Shows}
A horizontal bar chart showing spread over the last 10 time periods.

\subsubsection{Color Coding}
\begin{itemize}
    \item \textbf{Green:} Low spread ($<$3 bps) - Good liquidity
    \item \textbf{Yellow:} Medium spread (3-6 bps) - Normal
    \item \textbf{Red:} High spread ($>$6 bps) - Poor liquidity
\end{itemize}

\subsubsection{How to Interpret}
\begin{itemize}
    \item All green = excellent market conditions for trading
    \item Mostly red = be careful, high trading costs
    \item Transition from green to red = liquidity deteriorating
\end{itemize}

\subsection{Chart 3: Order Book Imbalance}

\subsubsection{What It Shows}
A gauge-style horizontal bar showing buy/sell pressure.

\subsubsection{Components}
\begin{itemize}
    \item \textbf{Green bar (left):} Bid volume (buyers)
    \item \textbf{Red bar (right):} Ask volume (sellers)
    \item \textbf{Percentage label:} The imbalance value
\end{itemize}

\subsubsection{How to Interpret}
\begin{itemize}
    \item Longer green bar = more buyers = bullish
    \item Longer red bar = more sellers = bearish
    \item Equal bars = balanced market
\end{itemize}

\subsection{Chart 4: Trade Velocity Gauge}

\subsubsection{What It Shows}
A speedometer-style gauge showing trading activity.

\subsubsection{Components}
\begin{itemize}
    \item \textbf{Needle:} Points to current velocity
    \item \textbf{Green zone (0-30):} Low activity
    \item \textbf{Yellow zone (30-60):} Normal activity
    \item \textbf{Red zone (60-100):} High activity / spike
    \item \textbf{White dashed line:} Baseline (average velocity)
\end{itemize}

\subsubsection{How to Interpret}
\begin{itemize}
    \item Needle in green = calm market
    \item Needle in red = unusual activity, possible news event
    \item Needle well above baseline = velocity spike, be alert
\end{itemize}

\subsection{Chart 5: Volatility Monitor}

\subsubsection{What It Shows}
An area chart showing rolling volatility in basis points (bps).

\subsubsection{Components}
\begin{itemize}
    \item \textbf{Purple area:} Volatility over time
    \item \textbf{Yellow dashed line at 15 bps:} Normal threshold
    \item \textbf{Red dashed line at 25 bps:} High volatility threshold
\end{itemize}

\subsubsection{How to Interpret}
\begin{itemize}
    \item Below yellow line = low risk, calm market
    \item Between yellow and red = moderate risk
    \item Above red line = HIGH RISK, reduce position sizes
\end{itemize}

\subsection{Chart 6: Trading Intelligence Panel}

This is the \textbf{KEY DIFFERENTIATOR} of AlgoViz!

\subsubsection{What It Shows}
AI-generated trading insights with:
\begin{itemize}
    \item \textbf{Priority badge:} 🔴 HIGH, 🟡 MEDIUM, 🟢 LOW
    \item \textbf{Insight:} What is happening in the market
    \item \textbf{Action:} What you should do (IF-THEN recommendation)
    \item \textbf{How to Overcome:} Specific implementation steps
    \item \textbf{Expected Impact:} Quantified benefit of following the advice
\end{itemize}

\subsubsection{Example Insight}
\begin{tcolorbox}[colback=red!10, colframe=red!50, title=🔴 HIGH Priority]
\textbf{Insight:} Spread widened to 7.5 bps — liquidity deteriorating\\
\textbf{Action:} Pause market orders; use limit orders only\\
\textbf{Overcome:} Set limit orders at mid-price; accept partial fills\\
\textbf{Impact:} Save 6-8 bps per trade
\end{tcolorbox}

%===============================================================================
% SECTION 8: LIVE FEED PAGE
%===============================================================================
\section{Live Feed Page (Live Mode Only)}

This page is \textbf{only available in Live Data mode} and shows raw streaming data.

\subsection{Tick-by-Tick Price Chart}
Shows the last 30 seconds of price data at individual trade (tick) level.

\subsection{Trade Stream Table}
Real-time table showing each trade:
\begin{itemize}
    \item \textbf{Time:} When the trade happened (HH:MM:SS.mmm)
    \item \textbf{Price:} Trade price
    \item \textbf{Quantity:} Amount of BTC traded
    \item \textbf{Side:} BUY (green) or SELL (red)
\end{itemize}

\subsection{Order Book Depth Visualization}
Two horizontal bar charts:
\begin{itemize}
    \item \textbf{Bid Depth (Green):} Shows 10 best bid levels with volumes
    \item \textbf{Ask Depth (Red):} Shows 10 best ask levels with volumes
\end{itemize}

\subsection{Stream Statistics}
Four cards showing:
\begin{itemize}
    \item \textbf{Total Trades:} Number of trades received
    \item \textbf{Bid Volume:} Total BTC on buy side
    \item \textbf{Ask Volume:} Total BTC on sell side
    \item \textbf{Latency:} Milliseconds since last data update
\end{itemize}

\subsection{Trade Activity Heatmap}
Stacked bar chart showing buy/sell trade counts over the last 60 seconds.

%===============================================================================
% SECTION 9: ANALYTICS PAGE
%===============================================================================
\section{Analytics Page: Advanced Analysis}

The Analytics page has \textbf{5-6 tabs} depending on mode.

\subsection{Tab 1: Distribution \& Returns}

\subsubsection{Price Distribution Chart}
\textbf{Type:} Histogram\\
\textbf{X-axis:} Price ranges\\
\textbf{Y-axis:} Frequency (count of prices in each range)\\
\textbf{Interpretation:} Shows where prices cluster. Tall bars = prices often at that level.

\subsubsection{Volatility Distribution Chart}
\textbf{Type:} Histogram\\
\textbf{X-axis:} Volatility ranges (bps)\\
\textbf{Y-axis:} Frequency\\
\textbf{Interpretation:} Shows typical volatility levels. Tall bars on left = usually calm market.

\subsubsection{Return Distribution Chart}
\textbf{Type:} Histogram\\
\textbf{X-axis:} Return percentages (can be negative)\\
\textbf{Y-axis:} Frequency\\
\textbf{Interpretation:} Bell-shaped = normal market. Fat tails = extreme moves more common.

\subsubsection{Rolling Volatility Chart}
\textbf{Type:} Area chart\\
\textbf{X-axis:} Time\\
\textbf{Y-axis:} Volatility percentage\\
\textbf{Interpretation:} Rising = increasing risk. Spikes often precede big price moves.

\subsubsection{Distribution Statistics}
Shows 6 metrics:
\begin{itemize}
    \item \textbf{Avg Price:} Mean price over the period
    \item \textbf{Price Std:} Standard deviation of prices
    \item \textbf{Min Price:} Lowest price seen
    \item \textbf{Max Price:} Highest price seen
    \item \textbf{Avg Volatility:} Mean volatility
    \item \textbf{Vol Std:} Volatility of volatility (meta-volatility)
\end{itemize}

\subsubsection{Return Statistics}
Shows 6 metrics:
\begin{itemize}
    \item \textbf{Mean Return:} Average return per period
    \item \textbf{Return Std:} Standard deviation of returns
    \item \textbf{Min Return:} Worst return (most negative)
    \item \textbf{Max Return:} Best return (most positive)
    \item \textbf{Skewness:} Asymmetry of returns (negative = more downside risk)
    \item \textbf{Kurtosis:} ``Fat tails'' - higher = more extreme events
\end{itemize}

\subsection{Tab 2: Relationships (Correlation)}

\subsubsection{Correlation Matrix}
Shows how different features move together:
\begin{itemize}
    \item \textbf{+1.0:} Perfect positive correlation (move together)
    \item \textbf{-1.0:} Perfect negative correlation (move opposite)
    \item \textbf{0.0:} No correlation (independent)
\end{itemize}

\subsubsection{Scatter Plots}
\begin{itemize}
    \item \textbf{Price Change vs Volume:} Tests if big moves have high volume
    \item \textbf{Momentum vs Imbalance:} Tests if order flow predicts price
\end{itemize}

Each has a trend line (OLS regression) showing the relationship direction.

\subsection{Tab 3: Signals \& Alerts}

\subsubsection{Trading Intelligence Panel}
Same as on Dashboard - shows AI-generated insights.

\subsubsection{Active Alerts Panel}
Shows triggered alerts with:
\begin{itemize}
    \item Priority level (Critical/High/Medium/Low/Info)
    \item Alert message
    \item Value that triggered the alert
    \item Timestamp
\end{itemize}

\subsection{Tab 4: Backtester}

\begin{importantbox}[\textbf{What is Backtesting?}]
Backtesting means testing a trading strategy on historical data to see how it would have performed. This helps validate strategies BEFORE risking real money.
\end{importantbox}

\subsubsection{Configuration Options}
\begin{itemize}
    \item \textbf{Select Strategy:} Choose from 4 pre-built strategies
    \item \textbf{Initial Capital:} Starting money (\$1,000 to \$1,000,000)
    \item \textbf{Simulation Periods:} How much data to test on (50-500)
\end{itemize}

\subsubsection{Available Strategies}

\begin{table}[h]
\centering
\begin{tabular}{lp{9cm}}
\toprule
\textbf{Strategy} & \textbf{How It Works} \\
\midrule
Momentum Breakout & Buys when imbalance $>$40\% AND velocity $>$1.5x baseline. Exits when imbalance reverses or after 30 seconds. \\
Mean Reversion & Buys when price is 0.15\% below VWAP AND volatility $<$20 bps. Sells when price returns to VWAP. \\
Volatility Breakout & Trades during high volatility ($>$22 bps) when there's strong imbalance ($>$50\%). Quick exits. \\
Spread Fade & Enters when spread is wide ($>$6 bps), expecting it to narrow. Exits when spread normalizes. \\
\bottomrule
\end{tabular}
\caption{Available Trading Strategies}
\end{table}

\subsubsection{Backtest Results}
After clicking ``Run Backtest'':

\textbf{Summary Cards:}
\begin{itemize}
    \item \textbf{Total P\&L:} Profit or loss in dollars
    \item \textbf{Win Rate:} Percentage of profitable trades
    \item \textbf{Total Trades:} Number of trades executed
    \item \textbf{Sharpe Ratio:} Risk-adjusted return (higher = better)
    \item \textbf{Max Drawdown:} Worst peak-to-trough decline
    \item \textbf{Profit Factor:} Gross profit / Gross loss
\end{itemize}

\textbf{Equity Curve Chart:}
\begin{itemize}
    \item \textbf{X-axis:} Time
    \item \textbf{Y-axis:} Portfolio value
    \item \textbf{Interpretation:} Upward slope = making money. Dips = drawdowns.
\end{itemize}

\textbf{Trade Distribution Chart:}
\begin{itemize}
    \item Histogram of P\&L per trade
    \item Bars on right = winning trades
    \item Bars on left = losing trades
\end{itemize}

\subsubsection{Model Evaluation}
\textbf{Confusion Matrix:}
Shows prediction accuracy:
\begin{itemize}
    \item Diagonal = correct predictions
    \item Off-diagonal = mistakes
    \item Labels: UP, HOLD, DOWN
\end{itemize}

\textbf{ROC Curves:}
\begin{itemize}
    \item Shows model discrimination ability per class
    \item AUC $>$0.9 = excellent
    \item AUC 0.7-0.9 = good to fair
    \item Closer to top-left = better performance
\end{itemize}

\subsection{Tab 5: Deep Learning}

This tab explains and demonstrates the neural network predictions.

\subsubsection{Deep Learning Panel}
Shows prediction results:
\begin{itemize}
    \item \textbf{Direction:} STRONG\_UP, UP, NEUTRAL, DOWN, STRONG\_DOWN
    \item \textbf{Confidence:} How sure the model is (0-100\%)
    \item \textbf{Predicted Move:} Expected price change in bps
    \item \textbf{Up/Down/Neutral Probabilities:} Model's probability distribution
    \item \textbf{Market Regime:} TRENDING\_UP, TRENDING\_DOWN, RANGING, VOLATILE, BREAKOUT
\end{itemize}

\subsubsection{Prediction Summary Card}
Quick view of:
\begin{itemize}
    \item \textbf{Recommended Action:} STRONG BUY, BUY, HOLD, SELL, STRONG SELL
    \item \textbf{Signal Strength:} 0-100\%
    \item \textbf{Prediction Horizon:} 5 seconds ahead
    \item \textbf{Expected Move:} In basis points
    \item \textbf{Model Confidence:} Uncertainty level
\end{itemize}

\subsubsection{Neural Network Architecture}
Visual breakdown of the model:

\textbf{Input Layer:}
\begin{itemize}
    \item Shape: (sequence\_length, 8 features)
    \item Features: Price, Momentum, Volatility, Imbalance, Spread, Volume, VWAP Deviation, Trade Velocity
\end{itemize}

\textbf{Hidden Layers:}
\begin{itemize}
    \item LSTM Cell (32 units) - captures temporal patterns
    \item Attention Layer - weights important timesteps
    \item Dense Layers (32 → 16 units) - final processing
\end{itemize}

\textbf{Output Heads:}
\begin{itemize}
    \item Direction Classifier: 3-class softmax (Up/Hold/Down)
    \item Regime Classifier: 5-class softmax
    \item Regression: Predicted move in bps
\end{itemize}

\subsubsection{Advanced Visualizations}
\begin{itemize}
    \item \textbf{Candlestick Chart:} Traditional OHLC (Open-High-Low-Close) visualization
    \item \textbf{Order Book Depth Chart:} Cumulative bid/ask volume
\end{itemize}

\subsubsection{Algorithm Comparison Table}
Compares 8 algorithms showing why LSTM+Attention is best for trading:
\begin{itemize}
    \item Linear/Logistic Regression - No temporal modeling
    \item Random Forest/XGBoost - Limited sequence memory
    \item SVM - No temporal modeling
    \item Simple RNN - Short-term memory only
    \item LSTM - Long-term memory ✓
    \item LSTM + Attention - Long-term memory + focus on important moments ✓✓
\end{itemize}

\subsection{Tab 6: Data Quality (Static Mode Only)}

Shows data cleaning and validation pipeline.

\subsubsection{Validation Summary}
\begin{itemize}
    \item \textbf{Total Rows:} All rows preserved (no data dropped!)
    \item \textbf{Values Imputed:} Missing values filled
    \item \textbf{Outliers Winsorized:} Extreme values capped
    \item \textbf{Duplicates Resolved:} New IDs assigned
\end{itemize}

\subsubsection{Cleaning Pipeline Steps}
\begin{enumerate}
    \item Missing Values → Forward-fill → Backward-fill → Median
    \item Negative Prices → Absolute value conversion
    \item Zero Quantities → Replace with median
    \item Duplicate IDs → Assign new unique IDs
    \item Timestamp Order → Sort \& interpolate
    \item Outliers → Winsorization (1st-99th percentile)
\end{enumerate}

\subsubsection{PCA Analysis}
\begin{itemize}
    \item \textbf{Scree Plot:} Shows variance explained by each component
    \item \textbf{Cumulative Line:} Shows how many components needed for 95\% variance
    \item \textbf{Dimensionality Reduction:} Reduces features while preserving information
\end{itemize}

%===============================================================================
% SECTION 10: TECHNICAL DEEP DIVE
%===============================================================================
\section{Technical Deep Dive}

\subsection{What is a WebSocket?}

\begin{importantbox}[\textbf{Simple Explanation}]
A WebSocket is like a phone call between your computer and Binance's server. Once connected, both sides can send messages instantly without hanging up. This is different from regular web requests (HTTP) which are like sending letters back and forth.
\end{importantbox}

\textbf{Why WebSocket for Trading?}
\begin{itemize}
    \item \textbf{Low Latency:} Data arrives in milliseconds
    \item \textbf{Real-time:} No need to keep asking ``any updates?''
    \item \textbf{Efficient:} Less network overhead than polling
\end{itemize}

\subsection{How Binance Data Comes to Dashboard}

\textbf{Step 1: Connection}
\begin{verbatim}
wss://stream.binance.com:9443/stream?streams=btcusdt@trade/btcusdt@depth10@100ms
\end{verbatim}
This URL connects to two streams simultaneously:
\begin{itemize}
    \item \texttt{btcusdt@trade} - Every single BTC/USDT trade
    \item \texttt{btcusdt@depth10@100ms} - Order book (10 levels) every 100ms
\end{itemize}

\textbf{Step 2: Data Parsing}
Each message is JSON containing:
\begin{itemize}
    \item For trades: price, quantity, timestamp, buy/sell side
    \item For depth: lists of bid prices/quantities, ask prices/quantities
\end{itemize}

\textbf{Step 3: State Management}
Data is stored in circular buffers:
\begin{itemize}
    \item Trades buffer: Last 1,000 trades
    \item Depth buffer: Last 100 order book snapshots
    \item Price history: Last 200 price points
\end{itemize}

\textbf{Step 4: Feature Calculation}
Every update, the Feature Engine calculates:
\begin{itemize}
    \item VWAP from last 30 seconds of trades
    \item Velocity from last 3 seconds
    \item Volatility from last 60 seconds
    \item Imbalance from latest order book
\end{itemize}

\textbf{Step 5: Display}
Streamlit fragments update every 1 second, showing the latest data.

\subsection{Is It Real-Time?}

\textbf{Yes, but with caveats:}
\begin{itemize}
    \item \textbf{Binance Server Time:} Data has Binance's server timestamp
    \item \textbf{Network Latency:} ~50-200ms delay due to internet
    \item \textbf{UI Refresh:} Dashboard updates every 500ms-1s
    \item \textbf{Overall Delay:} ~200-500ms from trade happening to seeing it
\end{itemize}

\textbf{Timezone:}
All timestamps are in \textbf{UTC (Coordinated Universal Time)}. The dashboard displays times in UTC format.

\subsection{Deep Learning: LSTM with Attention}

\subsubsection{What is LSTM?}

\begin{importantbox}[\textbf{Simple Explanation}]
LSTM (Long Short-Term Memory) is a type of neural network that can ``remember'' patterns over time. Unlike regular neural networks that only see one data point, LSTM processes sequences and remembers what came before.
\end{importantbox}

\textbf{Why LSTM for Trading?}
\begin{itemize}
    \item Markets have \textbf{temporal patterns} (e.g., momentum)
    \item Past prices \textbf{influence} future prices
    \item LSTM can learn these time-based relationships
\end{itemize}

\subsubsection{What is Attention?}

\begin{importantbox}[\textbf{Simple Explanation}]
Attention is like highlighting the important parts of a book. Instead of treating all past data equally, attention learns which past moments matter most for the current prediction.
\end{importantbox}

\textbf{Example:} If there was a big trade 5 seconds ago, attention might weight that moment heavily when predicting the next price move.

\subsubsection{Model Architecture}

\begin{enumerate}
    \item \textbf{Input:} 30 timesteps × 8 features = 240 inputs
    \item \textbf{LSTM Layer:} Processes sequence, outputs hidden states
    \item \textbf{Attention Layer:} Weights hidden states by importance
    \item \textbf{Dense Layers:} 32 → 16 neurons with ReLU activation
    \item \textbf{Output Heads:}
    \begin{itemize}
        \item Softmax for direction (3 classes: up/hold/down)
        \item Softmax for regime (5 classes)
        \item Linear for predicted move (regression)
    \end{itemize}
\end{enumerate}

\subsubsection{Features Used}
\begin{enumerate}
    \item \textbf{Price Return:} Percentage change from previous price
    \item \textbf{Momentum:} Rate of price change
    \item \textbf{Volatility:} Standard deviation of returns
    \item \textbf{Imbalance:} Order book imbalance
    \item \textbf{Spread:} Bid-ask spread
    \item \textbf{Volume:} Trading volume
    \item \textbf{VWAP Deviation:} Price vs VWAP percentage
    \item \textbf{Trade Velocity:} Trades per second
\end{enumerate}

%===============================================================================
% SECTION 11: PRESENTATION FLOW
%===============================================================================
\section{Recommended Presentation Flow}

\subsection{Suggested Order (15-20 minutes)}

\begin{enumerate}
    \item \textbf{Home Page (2 min)}
    \begin{itemize}
        \item Show animated landing page
        \item Highlight the 6 feature cards
        \item Mention team names
    \end{itemize}
    
    \item \textbf{Static Demo - Dashboard (4 min)}
    \begin{itemize}
        \item Switch to Static Demo mode
        \item Explain KPI cards (show the 8 metrics)
        \item Walk through each chart briefly
        \item Highlight the Trading Intelligence panel
    \end{itemize}
    
    \item \textbf{Static Demo - Analytics (3 min)}
    \begin{itemize}
        \item Show Distribution tab (histograms)
        \item Show Correlation tab
        \item Quick demo of Backtester (run one backtest)
    \end{itemize}
    
    \item \textbf{Deep Learning Tab (3 min)}
    \begin{itemize}
        \item Explain the prediction panel
        \item Show neural network architecture diagram
        \item Explain why LSTM+Attention (show comparison table)
    \end{itemize}
    
    \item \textbf{Live Demo (3 min)}
    \begin{itemize}
        \item Switch to Live Data mode
        \item Show real-time updates on Dashboard
        \item Show Live Feed page with streaming trades
        \item Show order book depth
    \end{itemize}
    
    \item \textbf{Technical Q\&A (2 min)}
    \begin{itemize}
        \item WebSocket explanation
        \item Docker setup
        \item Any technical questions
    \end{itemize}
\end{enumerate}

\subsection{Key Points to Emphasize}

\begin{enumerate}
    \item \textbf{Unique Value:} Trading Intelligence panel (not just charts!)
    \item \textbf{Real-Time:} Live Binance data with sub-second updates
    \item \textbf{AI/ML:} LSTM with Attention for predictions
    \item \textbf{Professional Quality:} Dark theme, glassmorphism, production-ready
    \item \textbf{Practical:} Backtester for strategy validation
    \item \textbf{Comprehensive:} 10+ chart types, 8 KPIs, 4 strategies
\end{enumerate}

%===============================================================================
% SECTION 12: COMMON QUESTIONS AND ANSWERS
%===============================================================================
\section{Common Questions and Answers}

\subsection{About the Data}

\textbf{Q: Where does the data come from?}\\
A: In Live mode, directly from Binance cryptocurrency exchange via WebSocket. In Static mode, it's synthetically generated by our algorithm.

\textbf{Q: Is this real money trading?}\\
A: No, this is a visualization and analysis dashboard. It does NOT execute real trades.

\textbf{Q: Why Bitcoin (BTC/USDT)?}\\
A: BTC/USDT is the most liquid cryptocurrency pair with high trading activity, making it ideal for demonstrating real-time capabilities.

\subsection{About the Technology}

\textbf{Q: Why Streamlit instead of React/Angular?}\\
A: Streamlit allows rapid development in Python, integrates seamlessly with data science libraries, and is perfect for data-focused dashboards.

\textbf{Q: Why not use TensorFlow/PyTorch for deep learning?}\\
A: We implemented LSTM in pure NumPy to minimize dependencies and demonstrate understanding of the algorithm's internals.

\textbf{Q: What is Docker for?}\\
A: Docker containerizes the application so it runs the same way on any computer, making deployment and demonstration reliable.

\subsection{About the Features}

\textbf{Q: What is a basis point (bps)?}\\
A: 1 basis point = 0.01\% = 0.0001. Used for precision in financial calculations.

\textbf{Q: How accurate is the deep learning prediction?}\\
A: The demo shows 54-72\% accuracy, which is realistic for short-term market prediction. Even small edges compound over many trades.

\textbf{Q: What does ``winsorization'' mean in data cleaning?}\\
A: Capping extreme values at the 1st and 99th percentiles instead of removing them. This preserves all data while reducing outlier impact.

%===============================================================================
% SECTION 13: TROUBLESHOOTING
%===============================================================================
\section{Troubleshooting During Demo}

\subsection{Dashboard Won't Load}
\begin{enumerate}
    \item Check if Docker container is running: \texttt{docker ps}
    \item Restart container: \texttt{docker-compose restart}
    \item Check port 8501 is not blocked
\end{enumerate}

\subsection{Live Data Not Updating}
\begin{enumerate}
    \item Check internet connection
    \item Look at sidebar debug info (WS Running: True?)
    \item Try refreshing the page (F5)
    \item Switch to Static mode and back to Live
\end{enumerate}

\subsection{Charts Show ``Waiting for data...''}
\begin{enumerate}
    \item Wait 5-10 seconds for data to accumulate
    \item Check connection status banner
    \item Ensure WebSocket is connected (green dot)
\end{enumerate}

\subsection{Backup Plan}
If live demo fails:
\begin{enumerate}
    \item Switch to Static Demo mode (works offline)
    \item All charts and features work with synthetic data
    \item Explain that live connection issues can occur but functionality is identical
\end{enumerate}

%===============================================================================
% SECTION 14: GLOSSARY
%===============================================================================
\section{Glossary of Terms}

\begin{longtable}{lp{10cm}}
\toprule
\textbf{Term} & \textbf{Definition} \\
\midrule
\endhead
Ask & The lowest price a seller is willing to accept \\
Bid & The highest price a buyer is willing to pay \\
BPS (Basis Points) & One hundredth of a percent (0.01\%) \\
Drawdown & Peak-to-trough decline in portfolio value \\
Imbalance & Difference between buy and sell volume \\
LSTM & Long Short-Term Memory (neural network type) \\
Market Order & Order to buy/sell immediately at current price \\
Order Book & List of all open buy and sell orders \\
P\&L & Profit and Loss \\
ROC Curve & Receiver Operating Characteristic - model evaluation metric \\
Sharpe Ratio & Return per unit of risk \\
Spread & Difference between bid and ask prices \\
TWAP & Time-Weighted Average Price \\
Velocity & Number of trades per second \\
Volatility & Measure of price variability \\
VWAP & Volume-Weighted Average Price \\
WebSocket & Real-time bidirectional communication protocol \\
Winsorization & Capping extreme values at percentiles \\
\bottomrule
\caption{Glossary of Trading and Technical Terms}
\end{longtable}

%===============================================================================
% SECTION 15: CONCLUSION
%===============================================================================
\section{Conclusion}

\begin{importantbox}[\textbf{Summary}]
AlgoViz is a comprehensive, professional-grade algorithmic trading visualization dashboard that combines:
\begin{itemize}
    \item Real-time data streaming from Binance
    \item Deep learning predictions using LSTM with Attention
    \item 10+ interactive chart types
    \item AI-powered trading insights
    \item Strategy backtesting capabilities
    \item Professional dark theme design
\end{itemize}
\end{importantbox}

\subsection{What Makes AlgoViz Stand Out}
\begin{enumerate}
    \item \textbf{Not Just Charts:} The Trading Intelligence panel provides actionable insights
    \item \textbf{Real AI:} Genuine LSTM neural network implementation
    \item \textbf{Production Quality:} Docker deployment, proper error handling
    \item \textbf{Educational:} Great for learning algorithmic trading concepts
    \item \textbf{Practical:} Backtester validates strategies before trading
\end{enumerate}

\vspace{1cm}

\begin{center}
\textbf{\Large Good luck with your presentation!}

\vspace{0.5cm}

\textit{``See the Signal in the Noise''}

\vspace{1cm}

\textcolor{algovizpurple}{\textbf{AlgoViz Team}}\\
Kartik Joshi $\bullet$ Aditya Chitale $\bullet$ Krishna Patel
\end{center}

\end{document}
